% 演示文稿：面向 CODESYS 平台的 ST 代码自动生成系统
% 编译：xelatex 两次（需要 texlive，使用 XeLaTeX）
\PassOptionsToPackage{quiet}{fontspec}
\PassOptionsToPackage{dvipsnames}{xcolor}
\documentclass[10pt,aspectratio=169]{beamer}
\usepackage{zju_beamer}
\zjusectioncolor{blue}
\usepackage{colortbl}
\usepackage{tikz}
\usetikzlibrary{arrows.meta,positioning,shapes.geometric}

\graphicspath{{./figures/}}

% 首页信息设置
\title[CODESYS ST 代码自动生成]{面向 CODESYS 平台的 ST 代码自动生成系统}
\subtitle{——自研 Agent Harness 与现有框架的四设置对比}
\author[报告人]{报告人}
\institute[ZJU]{计算机学院 \\ 浙江大学}
\date{\today}

% 顶部栏隐藏导航，节省空间（可选）
\setbeamertemplate{section in head/foot shaded}[default]

\begin{document}

% ---------------------------------------------------------------- 1 标题页
\begin{frame}
	\titlepage
\end{frame}

% ---------------------------------------------------------------- 2 任务与难点
\begin{frame}{问题：让大模型写出"真的能跑"的 PLC 代码}
	\begin{columns}[T]
		\begin{column}{0.52\textwidth}
			\begin{block}{任务}
				输入自然语言需求 + 接口表 + LTL 规约 \\
				输出可直接在 CODESYS 编译并运行的
				\textbf{IEC 61131-3 结构化文本（ST）}
			\end{block}
			\begin{alertblock}{三个难点}
				\begin{itemize}
					\item 语法/工程结构：DUT、Standard 库、调用图
					\item 行为语义：状态机、定时器、边沿、优先级
					\item 评测可信：需要\textbf{真实执行}而非文本相似
				\end{itemize}
			\end{alertblock}
		\end{column}
		\begin{column}{0.46\textwidth}
			\begin{block}{为什么不能只看"像不像"}
				\begin{itemize}
					\item 编译通过 $\neq$ 行为正确：
					      差一拍、优先级错都会"看起来对"
					\item 自研解释器又会引入
					      "不支持某语法"的假失败
				\end{itemize}
			\end{block}
			\begin{exampleblock}{判定标准}
				必须来自\textbf{真实 CODESYS 编译 + 仿真运行}，
				而不是文本相似度或自研解释器
			\end{exampleblock}
		\end{column}
	\end{columns}
\end{frame}

% ---------------------------------------------------------------- 3 数据集
\begin{frame}{数据集：可编译、可执行、可追溯}
	\begin{columns}[T]
		\begin{column}{0.46\textwidth}
			\begin{block}{规模}
				\begin{itemize}
					\item \textbf{30 题}：公共 20 + 自建 10
					\item 难度：easy 9 / medium 12 / hard 9
					\item \textbf{75 条 LTL 规约}（2.5 条/题）
					\item \textbf{54 组用例 / 199 个扫描步}
				\end{itemize}
			\end{block}
			\begin{block}{每题包含}
				中英需求 $\cdot$ 结构化接口（含 DUT）$\cdot$
				LTL 规约 $\cdot$ 行为用例 $\cdot$
				\textbf{可编译的参考实现}
			\end{block}
		\end{column}
		\begin{column}{0.5\textwidth}
			\begin{exampleblock}{质量门槛：先证明"题可做、用例对"}
				\begin{itemize}
					\item 参考实现 \textbf{30/30} 编译通过
					\item 参考实现 \textbf{30/30} 通过自己的用例
					\item 首轮自检修正 \textbf{16 处}
					      用例/参考实现缺陷
				\end{itemize}
			\end{exampleblock}
			\begin{block}{用例覆盖}
				正常流程 $\cdot$ 边界条件 $\cdot$
				优先级冲突 $\cdot$ 错误处理；
				\textbf{每步写明测试目的}，并关联到 LTL 规约
			\end{block}
		\end{column}
	\end{columns}
\end{frame}

% ---------------------------------------------------------------- 4 系统架构
\begin{frame}{系统架构：编译与仿真两道闸门 + 结构化反馈}
	\centering
	\resizebox{0.97\textwidth}{!}{%
\begin{tikzpicture}[
			x=1mm,y=1mm,
			box/.style={draw=blue!45,fill=blue!6,rounded corners,
					align=center,font=\scriptsize,inner sep=2pt,
					minimum height=9mm,text width=21mm},
			gate/.style={draw=orange!75,fill=orange!12,rounded corners,
					align=center,font=\scriptsize,inner sep=2pt,
					minimum height=9mm,text width=21mm},
			ok/.style={draw=green!45!black,fill=green!10,rounded corners,
					align=center,font=\scriptsize,inner sep=2pt,
					minimum height=9mm,text width=21mm},
			ar/.style={-{Latex[length=1.5mm]},thick,blue!55!black},
			back/.style={-{Latex[length=1.5mm]},thick,orange!80!black,dashed}]

		\node[box]  (task) at (11,0)  {题目\\需求+接口+LTL};
		\node[box]  (llm)  at (37,0)  {LLM\\生成整块代码};
		\node[box]  (chk)  at (63,0)  {解析+接口校验\\零成本};
		\node[gate] (cmp)  at (89,0)  {编译\\（CODESYS）};
		\node[gate] (sim)  at (115,0) {行为测试\\（CODESYS 仿真）};
		\node[ok]   (done) at (139,0) {成功};

		\draw[ar] (task) -- (llm);
		\draw[ar] (llm) -- (chk);
		\draw[ar] (chk) -- (cmp);
		\draw[ar] (cmp) -- (sim);
		\draw[ar] (sim) -- (done);

		\node[fill=orange!8,draw=orange!50,rounded corners,font=\scriptsize,
		      inner sep=2.5pt,text width=118mm,align=center] (fb) at (70,-17)
		     {结构化反馈：接口/命名错误 $\cdot$ 编译错误码+行列+消息
		      $\cdot$ 失败用例/步+期望值+实测值};
		\draw[back] (chk.south) -- (chk.south |- fb.north);
		\draw[back] (cmp.south) -- (cmp.south |- fb.north);
		\draw[back] (sim.south) -- (sim.south |- fb.north);
		\draw[back] (fb.west) -- (fb.west -| llm.south) -- (llm.south);
	\end{tikzpicture}
	}

	\vspace{0.4em}
	\begin{columns}[T]
		\begin{column}{0.49\textwidth}
			\begin{block}{上下文与成本控制}
				\begin{itemize}
					\item 只保留最近 2 轮对话
					\item 编译错误按 (code, text) 去重计数
					\item 测试只回传首个失败步
				\end{itemize}
			\end{block}
		\end{column}
		\begin{column}{0.47\textwidth}
			\begin{block}{终止条件}
				\begin{itemize}
					\item 通过 / 5 轮 / 代码无变化
					\item 失败签名连续 2 次相同
					\item \texttt{tool\_error} 单列，不计入模型
				\end{itemize}
			\end{block}
		\end{column}
	\end{columns}
\end{frame}

% ---------------------------------------------------------------- 5 CODESYS 接入
\begin{frame}{CODESYS 接入：用官方运行时，而不是自研解释器}
	\begin{columns}[T]
		\begin{column}{0.5\textwidth}
			\begin{block}{\texttt{codesys/compile}：编不编得过}
				\begin{itemize}
					\item 模板工程 + DUT + Standard 库注入
					\item \texttt{auto\_wire}：生成隐藏调用者，
					      让每个 POU 真正进入编译图
					\item 返回错误码 + 行列 + 中文消息
					\item 单次约 \textbf{13 s}
				\end{itemize}
			\end{block}
		\end{column}
		\begin{column}{0.48\textwidth}
			\begin{block}{\texttt{codesys/simulate}：跑起来对不对}
				\begin{itemize}
					\item 用例编译成 PLC 内\textbf{自测驱动}
					\item 逐拍设输入、断言输出、记录实测值
					\item 语义与编译器一致，无"不支持语法"
					\item 单次约 \textbf{24 s}
				\end{itemize}
			\end{block}
		\end{column}
	\end{columns}
	\vspace{0.4em}
	\begin{alertblock}{设计取舍}
		编译闸门（实测约 13 s）比仿真闸门（约 24 s）便宜一半，
		且格式/接口错误在\textbf{零成本预校验}就被拦下；
		对弱模型，多数修复轮次消耗在编译反馈上
		（Codex+Qwen 首轮 25/30 未编译通过）。
	\end{alertblock}
\end{frame}

% ---------------------------------------------------------------- 6 实验设置
\begin{frame}{实验设置：2×2 公平对比}
	\begin{columns}[T]
		\begin{column}{0.38\textwidth}
			\centering
			\renewcommand{\arraystretch}{1.3}
			\resizebox{\columnwidth}{!}{%
			\begin{tabular}{lcc}
				\hline
				& \textbf{DeepSeek} & \textbf{Qwen} \\
				\hline
				\textbf{Codex CLI} & 实验 1 & 实验 2 \\
				\textbf{自研 harness} & 实验 3 & 实验 4 \\
				\hline
			\end{tabular}
			}
			\vspace{0.8em}
			\begin{block}{规模}
				30 题 × 4 设置 = \textbf{120 次运行} \\
				每题最多 5 轮；\texttt{tool\_error}
				已重试至 \textbf{0}
			\end{block}
		\end{column}
		\begin{column}{0.57\textwidth}
			\begin{block}{公平性约束}
				\begin{itemize}
					\item 同一 DMXAPI 网关、同一模型快照
					\item 同一提示、同一反馈文本、同一轮次上限
					\item 同一 CODESYS 模板与 worker、
					      同一编译/测试管线
					\item Codex 包装成相同的
					      \texttt{chat(messages)} 接口，
					      差异只在生成组织方式
				\end{itemize}
			\end{block}
		\end{column}
	\end{columns}
\end{frame}

% ---------------------------------------------------------------- 7 结果
\begin{frame}{结果：DeepSeek 下持平，Qwen 下自研显著更好}
	\vspace{0.6em}
	\begin{center}
		\large
		\setlength{\tabcolsep}{7pt}
		\renewcommand{\arraystretch}{1.35}
		\begin{tabular}{llcccc}
			\hline
			框架 & 模型 & 首次编译 & 测试通过 & 平均轮次 & 输入 tokens \\
			\hline
			自研 & DeepSeek & 96.7\% & \textbf{83.3\%} & 1.77 & 75k \\
			自研 & Qwen     & 76.7\% & \textbf{66.7\%} & 2.53 & 131k \\
			Codex & DeepSeek & 100.0\% & \textbf{83.3\%} & 1.70 & 565k \\
			\rowcolor{red!12}
			Codex & Qwen     & 16.7\% & \textbf{33.3\%} & 4.00 & 1633k \\
			\hline
		\end{tabular}
	\end{center}
	\vspace{0.4em}
	\begin{columns}[T]
		\begin{column}{0.32\textwidth}
			\begin{block}{质量}
				DeepSeek 两框架测试通过率\textbf{持平}（83.3\%），
				自研最终编译 100\%
			\end{block}
		\end{column}
		\begin{column}{0.32\textwidth}
			\begin{block}{成本}
				自研输入 token 仅为 Codex 的 \textbf{13\%}
				（75k vs 565k），耗时约 \textbf{48\%}
			\end{block}
		\end{column}
		\begin{column}{0.32\textwidth}
			\begin{block}{鲁棒性}
				Qwen 经 Codex 首轮编译仅 16.7\%，
				18/30 题耗尽 5 轮
			\end{block}
		\end{column}
	\end{columns}
\end{frame}

% ---------------------------------------------------------------- 8 分析
\begin{frame}{分析：框架×模型存在强交互}
	\begin{columns}[T]
		\begin{column}{0.52\textwidth}
			\begin{block}{为什么 Qwen 经 Codex 会崩}
				\begin{itemize}
					\item 首轮普遍非 IEC 语法：
					      \texttt{C6} 140 次（把 \texttt{:=} 写成 \texttt{=}、用 \texttt{!}）
					\item 遗漏 \texttt{VAR\_IN\_OUT} 参数、非法 DUT
					\item 每轮修复成功率低 → 轮次耗尽
				\end{itemize}
			\end{block}
			\begin{block}{自研为何更稳}
				\begin{itemize}
					\item 块格式 + 接口预校验把错误挡在编译前
					\item 反馈短、聚焦单点，弱模型也能迭代
				\end{itemize}
			\end{block}
		\end{column}
		\begin{column}{0.46\textwidth}
			\begin{block}{剩余失败集中在行为语义}
				\begin{itemize}
					\item 状态机（五工位、柔性加工）
					\item 定时/边沿（焊缝去抖、自动门、交通灯）
					\item 典型症状：差一拍、优先级未覆盖
				\end{itemize}
			\end{block}
			\begin{alertblock}{结论}
				框架与模型需要\textbf{匹配调优}；
				"更强的 agent 框架对任意模型都更好"不成立。
			\end{alertblock}
		\end{column}
	\end{columns}
\end{frame}

% ---------------------------------------------------------------- 9 结论与局限
\begin{frame}{总结与局限}
	\begin{columns}[T]
		\begin{column}{0.5\textwidth}
			\begin{block}{已完成}
				\begin{itemize}
					\item 30 题数据集：LTL 规约 + 行为用例 + 参考实现
					\item 自研 harness：编译/仿真双闸门闭环
					\item CODESYS 编译与仿真服务，可整体迁移
					\item 四设置公平评测，自动测试（加分项）
				\end{itemize}
			\end{block}
		\end{column}
		\begin{column}{0.48\textwidth}
			\begin{alertblock}{局限与后续}
				\begin{itemize}
					\item 每题仅采样 1 次，存在方差
					\item LTL 未机器验证；用例覆盖有缺口
					\item 下一步：接入 PLCVerif、
					      补强用例覆盖、多平台验证
				\end{itemize}
			\end{alertblock}
		\end{column}
	\end{columns}
	\vspace{0.6em}
	\begin{center}
		\begin{tcolorbox}[colback=blue!6,colframe=blue!45,arc=2pt,
				left=4pt,right=4pt,top=2pt,bottom=2pt]
			\centering\small
			一句话结论：\textbf{把"能不能编译、跑得对不对"做成结构化反馈}，
			弱模型也能被推上来；框架与模型需要匹配调优。
		\end{tcolorbox}
	\end{center}
	\vspace{0.3em}
	\centering
	{\Large 谢谢，请指正}
\end{frame}

\end{document}
